\section{Introduction}

\lettrine{T}{echnology} scaling beyond 5\,nm has made power delivery harder. Transistor dimensions keep shrinking, but interconnects have not scaled at the same rate, which increases wire resistance and tightens voltage margins. Conventional front-side power delivery networks (PDNs) share routing resources with signal interconnects, and at these nodes they struggle to meet power integrity targets without giving up routing tracks and, in turn, performance~\cite{8993617}.

To work around these limitations, buried power rails (BPR) were proposed as an intermediate step. BPR places the power rails below the active devices using high-aspect-ratio ruthenium structures, which lowers local resistance and frees front-side metal layers for signal routing~\cite{8993617}. Backside power delivery networks (BS-PDN) go further: they move the entire power grid to the wafer backside, separating power and signal routing altogether~\cite{10185208}. Intel's PowerVia demonstration showed a 30\,\% reduction in voltage droop and a 6\,\% frequency gain on a fabricated test vehicle~\cite{10185208}. Separately, block-level studies at imec reported area savings and better routing efficiency in both high-performance and high-density standard-cell libraries~\cite{10631463}, and modelling work estimated more than a four-fold reduction in IR-drop relative to a conventional front-side PDN~\cite{8932458}.

These results point to a shift in how power networks are designed, from incremental back-end-of-line changes to a full reorganisation of the power distribution path. This literature review examines that shift: the move from front-side PDNs through buried power rails to backside delivery, weighing the reported power integrity gains against the new design challenges that these architectures introduce in sub-5\,nm CMOS technologies.
